\section{Zero modes of the Standard Model}\label{rep:zeromodes}

The constant modes of the gauge and Higgs fields are not free fields, and their energy is not part of the image sum of Sec.~\ref{rep:casimir}.
We show that their ground energy $Q/L$ is positive and independent of the shape at fixed volume, and that it uses up a fraction
\begin{equation}
\fzm\equiv Q/C=0.30 %
\label{eq:rA-fzero}
\end{equation}
of the free Casimir energy, between the lower bound $0.26$ and the Gaussian variational bound $0.31$. %
The Higgs zero mode is held mainly by the zero-point fluctuations of the electroweak gauge modes.
At next order in the small-volume expansion the fraction falls to about $0.28$ and the zero modes acquire a weak shape dependence.
Sections~\ref{rA:zm-hamiltonian}--\ref{rA:zm-isospectral} derive the Hamiltonian and its shape independence, Secs.~\ref{rA:zm-bounds}--\ref{rA:zm-couplings} its ground energy, Secs.~\ref{rA:zm-nextorder}--\ref{rA:zm-shape} the next order, and Sec.~\ref{rA:zm-consequences} the consequences for the background and for the Higgs zero mode.

\subsection{The constant modes and their Hamiltonian}
\label{rA:zm-hamiltonian}

Each of the twelve gauge fields has three constant modes, one per cycle, and the Higgs doublet has four.
Their state is not fixed by the Hadamard condition~\cite{Negro:2026ooh}, which constrains only the short-distance behavior, and their nonlinear self-interaction gives their ground state a finite energy~\cite{Luscher:1982uv,Luscher:1982ma}.
For antiperiodic fermions the vacuum lies at trivial holonomy and is unique~\cite{vanBaal:1988va}, and the AAA fermions themselves have no constant modes.

\paragraph{Reduction of the action.}
We use the covariant derivative $D=\partial-i(A^aT_a+YB)$ with ${\rm tr}\,T_aT_b=\tfrac12\delta_{ab}$, temporal gauge, and constant modes
\begin{equation}
A^a_i=g\,x^a_i(t),\qquad B_i=g_Y\,w_i(t)
\label{eq:rA-xw}
\end{equation}
in coordinates of unit period, with $g=g_3$ or $g_2$ for the non-abelian factors.
For constant spatial components the field strength is $F^a_{0i}=\dot A^a_i$ and $F^a_{ij}=f^{abc}A^b_iA^c_j=g^2F^a_{ij}[x]$, with $F^a_{ij}[x]=f^{abc}x^b_ix^c_j$.
With $\sqrt{-g}=L^3$ and $g^{ij}=G^{ij}/L^2$ the Yang--Mills action $-\frac{1}{4g^2}\int\sqrt{-g}\,F^a_{\mu\nu}F^{a\mu\nu}$ of each gauge factor reduces to
\begin{equation}
L_{\rm YM}=\frac{L}{2}\,G^{ij}\dot x_i\cdot\dot x_j-\frac{g^2}{4L}\,G^{ik}G^{jl}F_{ij}[x]\cdot F_{kl}[x],
\end{equation}
where the dot product sums over the adjoint index.
The constant Higgs doublet enters as $\Phi/L$, with $\Phi=(\phi_0+i\phi_1,\phi_2+i\phi_3)/\sqrt2$ and $r=|\bm\phi|$, so that $\Phi^\dagger\Phi=r^2/2$.
At fixed volume its kinetic term is $L\,\dot\Phi^\dagger\dot\Phi=\tfrac{L}{2}\sum_a\dot\phi_a^2$.
Its gradient term, $-\sqrt{-g}\,g^{ij}(D_i\Phi)^\dagger D_j\Phi$ with $D_i\Phi=-iM_i\Phi/L$ and
\begin{equation}
M_i=g_2\,x_i\cdot T+g_YY\,w_i ,
\end{equation}
contributes the potential $G^{ij}\Phi^\dagger M_iM_j\Phi/L$, and the quartic contributes $\lambda(\Phi^\dagger\Phi)^2/L$.
The Legendre transform, with momenta $\pi_i=LG^{ij}\dot x_j$, $\pi^w_i=LG^{ij}\dot w_j$ and $p_a=L\dot\phi_a$, gives the zero-mode Hamiltonian
\begin{equation}
\begin{aligned}
L\mathcal H_0={}&\sum_{\rm G}\Big[\tfrac12G_{ij}\,\pi_i\cdot\pi_j+\frac{g^2}{4}G^{ik}G^{jl}F_{ij}\cdot F_{kl}\Big]\\
&+\tfrac12G_{ij}\,\pi^w_i\pi^w_j+\tfrac12p_\Phi^2\\
&+G^{ij}\Phi^\dagger M_iM_j\Phi+\lambda(\Phi^\dagger\Phi)^2,
\end{aligned}
\label{eq:rA-H0}
\end{equation}
where the sum runs over the gauge factors $SU(3)$ and $SU(2)$, $F^a_{ij}=f^{abc}x^b_ix^c_j$, and $p_\Phi^2$ sums over the four real components $\phi_a$.
The Gauss law, $f^{abc}x^b_i\pi^c_i$ for each non-abelian factor plus the Higgs charge density, annihilates physical states.
The term $\lambda v^2\Phi^\dagger\Phi$ of the potential is negligible at $\mub$ and is omitted.

\subsection{Isospectrality}
\label{rA:zm-isospectral}

\begin{rAprop}[Shape independence]\label{prop:rA-isospectral}
On the local domain of the constant modes, the Hamiltonian~\eqref{eq:rA-H0} at shape $G$ is unitarily equivalent to the one at the cube.
Its whole local spectrum, not only the ground energy, is therefore independent of the shape at fixed volume.
\end{rAprop}

\begin{rAproof}
Write $G=SS^{\mathsf T}$ with $\det S=1$ and define the canonical map
\begin{equation}
x_i=S_{ij}\tilde x_j,\qquad w_i=S_{ij}\tilde w_j,\qquad \pi=S^{-\mathsf T}\tilde\pi,\qquad \pi^w=S^{-\mathsf T}\tilde\pi^w,
\label{eq:rA-isomap}
\end{equation}
with $\Phi$ unchanged.
The kinetic terms become $\pi^{\mathsf T}G\pi=\tilde\pi^{\mathsf T}S^{-1}SS^{\mathsf T}S^{-\mathsf T}\tilde\pi=\tilde\pi^2$.
The field strength is bilinear, $F_{ij}[x]=S_{ia}S_{jb}F_{ab}[\tilde x]$, and $S^{\mathsf T}G^{-1}S=\rAone$, so $G^{ik}G^{jl}F_{ij}\cdot F_{kl}=F[\tilde x]^2$.
Since $M_i$ is linear in $(x,w)$, $M_i(x,w)=S_{ij}M_j(\tilde x,\tilde w)$, and the same identity gives $G^{ij}\Phi^\dagger M_iM_j\Phi=\sum_j|M_j(\tilde x,\tilde w)\Phi|^2$.
The Gauss generators are unchanged, $\sum_ix_i^b\pi_i^c=\tilde x^b\cdot\tilde\pi^c$, and the Jacobian of the configuration-space map is $\det(\rAone\otimes S)\det S=1$, so the flat Laplacian of the kinetic term is preserved.
The map is therefore unitary on the local domain and carries $\mathcal H_0(G)$ into $\mathcal H_0(\rAone)$.
\end{rAproof}

We checked this exactly for a rational $S$ and numerically for 200 random $S$, to $2\times10^{-14}$. %
The ground energy $Q/L$ is therefore a positive constant on the space of shapes at fixed volume.
Only global features depend on the shape: the lattice of large holonomies is not invariant under the map, so their identification, tunneling between them and the matching to the nonzero modes can depend on the shape; for AAA fermions the ground state sits at trivial holonomy~\cite{vanBaal:1988va}, far from the copies (Sec.~\ref{rA:zm-consequences}).

\subsection{Small-volume scaling and bounds}
\label{rA:zm-bounds}

\paragraph{Scaling.}
At the cube the rescaling $x=g^{-1/3}y$, $\pi_x=g^{1/3}\pi_y$ turns the gauge part of each factor into
\begin{equation}
\tfrac12\pi_x^2+\frac{g^2}{4}F[x]^2=g^{2/3}\Big(\tfrac12\pi_y^2+\tfrac14F[y]^2\Big),
\end{equation}
whose ground energy defines $\epsilon_N$ for $SU(N)$.
At leading order of the small-volume expansion the pure-gauge part of $Q$ is therefore $\epsilon_3g_3^{2/3}+\epsilon_2g_2^{2/3}$, with $\epsilon_2=4.1167$~\cite{Luscher:1983gm} and $\epsilon_3=12.587$~\cite{Weisz:1986vu,Pavel:2021mxn}, independent of the fermions at this order~\cite{vanBaal:1988va}. %
With $g_3\simeq0.54$, $g_2\simeq0.53$ and $g_Y\simeq0.44$ at $\mub$, this is $26\%$ of the free Casimir coefficient $C=\NC\rAgam=42.74$. %

\paragraph{Lower bound.}
Write the electroweak part as $\mathcal H_{\rm EW}=\mathcal H_{SU(2)}\otimes\rAone+K$, where $K$ contains the hypercharge and Higgs kinetic terms, the Higgs gradient term and the quartic.
For $\lambda\geq0$ every term of $K$ is non-negative, $K\geq0$, so ${\rm inf\,spec}\,\mathcal H_{\rm EW}\geq\epsilon_2g_2^{2/3}/L$.
Restricting to gauge-invariant states can only raise the infimum, so
\begin{equation}
Q\geq\epsilon_2g_2^{2/3}+\epsilon_3g_3^{2/3}\qquad(\lambda\geq0).
\end{equation}
Because $\epsilon_3$ is a variational value~\cite{Weisz:1986vu,Pavel:2021mxn}, confirmed by the Monte Carlo below to $2\times10^{-4}$, this bound is numerical rather than rigorous. %
For $\lambda\geq0$ the ground state exists because the zero-point energy lifts the flat valleys of the potential, and it is unique and positive because $e^{-t\mathcal H_0}$ is positivity improving on the connected configuration space; it is therefore gauge invariant.
For $\lambda<0$ the Higgs direction is unbounded below and the ground state becomes a resonance, which Secs.~\ref{rep:h-matching} and~\ref{rep:h-reduced} treat.

\paragraph{Gaussian upper bound.}
The Gaussian
\begin{equation}
\psi\propto\exp\big[-\tfrac12\big(\alpha\,x^2+\beta\,w^2+z\,r^2\big)\big]
\end{equation}
for the $SU(2)$, hypercharge and Higgs zero modes has $\langle(x^a_i)^2\rangle=1/2\alpha$ for each of the nine $SU(2)$ components, $\langle w_i^2\rangle=1/2\beta$ for each of the three hypercharge components and $\langle\phi_a^2\rangle=1/2z$ for each of the four Higgs components.
Its energy follows term by term.
The kinetic terms give $\tfrac12\cdot9\cdot\tfrac\alpha2+\tfrac12\cdot3\cdot\tfrac\beta2+\tfrac12\cdot4\cdot\tfrac z2$.
The magnetic term gives $\tfrac{g_2^2}{4}\langle F\cdot F\rangle$ with $\langle F^a_{ij}F^a_{ij}\rangle=\sum_{i\neq j}\epsilon^{abc}\epsilon^{abc}/(2\alpha)^2=6\cdot6/4\alpha^2$.
The Higgs gradient term gives $g_2^2\sum_i\langle(x_i\cdot T)^2\rangle\langle\Phi^\dagger\Phi\rangle=g_2^2\cdot3\cdot\tfrac34\cdot\tfrac1{2\alpha}\cdot\tfrac1z$, using the doublet Casimir $T^aT^a=\tfrac34$ and $\langle\Phi^\dagger\Phi\rangle=\tfrac12\langle r^2\rangle=1/z$, and the hypercharge term the analogous $g_Y^2Y_H^2\cdot\tfrac3{2\beta}\cdot\tfrac1z$ with $Y_H=\tfrac12$.
The quartic gives $\tfrac\lambda4\langle r^4\rangle=\tfrac\lambda4\cdot\tfrac{6}{z^2}$.
Altogether the upper bound is
\begin{equation}
LE=\tfrac94\alpha+\tfrac34\beta+z+\frac{9g_2^2}{4\alpha^2}+\frac{9g_2^2}{8\alpha z}+\frac{3g_Y^2}{8\beta z}+\frac{3\lambda}{2z^2},
\label{eq:rA-gauss-bound}
\end{equation}
to be minimized over $\alpha$, $\beta$ and $z$.
For $SU(N)$ alone the Gaussian energy is $\tfrac34(N^2-1)\alpha+\tfrac38N(N^2-1)g^2/\alpha^2$; its minimum, at $\alpha^3=Ng^2$, is $\tfrac98(N^2-1)N^{1/3}g^{2/3}$, that is $4.252$ and $12.980$ in units of $g^{2/3}$ for $N=2$ and $3$, $3\%$ above $\epsilon_2$ and $\epsilon_3$. %

\subsection{Diffusion Monte Carlo}
\label{rA:zm-dmc}

For $\lambda\geq0$ the ground state has no nodes, so importance-sampled diffusion Monte Carlo (DMC), guided by the optimal Gaussian, gives the exact local ground energy up to statistical and time-step errors.
Walkers diffuse in the sixteen-dimensional configuration space of the $SU(2)$, hypercharge and Higgs modes with the drift of the guiding Gaussian and branch with the local energy.
We used $10^4$ walkers and time steps $0.04$, $0.02$ and $0.01$, down to $0.005$ for the checks, with linear extrapolation to zero step. %
The method reproduces the certified ground energy $2.1407$ of the four-dimensional quartic oscillator $\tfrac12p^2+\tfrac14(r^2)^2$ as $2.1402(5)$, $\epsilon_2=4.116720$ as $4.1171(14)$ and $\epsilon_3=12.5868$ as $12.5869(22)$. %
With the couplings of Sec.~\ref{rA:zm-couplings}, the coupled electroweak sector has $LE_{\rm EW}=4.4478(13)$ at $\lambda=0$ and $4.4750(17)$ at $\lambda=0.01$, so $Q=12.790$ and $12.818$, and $\fzm=0.2993$ and $0.2999$. %
At $\lambda=0.1$ the fraction is $0.3043$, and $g_2=g_3=0.5$, $g_Y=0.46$ give $0.2878$ at $\lambda=0.01$. %
The pure-gauge bound gives $\fzm\geq0.2580$ and the Gaussian bound $\fzm\leq0.3083$ at $\lambda=0.01$. %
A diffusion Monte Carlo computation of the coupled $SU(2)\times U(1)\times$Higgs ground state therefore gives $Q=12.8$ and Eq.~\eqref{eq:rA-fzero}. %

\paragraph{The Higgs zero mode is held by the gauge modes.}
At $\lambda=0.01$ the hypercharge and Higgs modes add $1.78/L$ to the $SU(2)$ energy, $3.9$ times the energy $2.1407\lambda^{1/3}/L=0.46/L$ of a decoupled quartic oscillator $\tfrac12p_\Phi^2+\lambda(\Phi^\dagger\Phi)^2$. %
Together with hypercharge, the Higgs zero mode raises the $SU(2)$ energy by about four times what a decoupled quartic oscillator would contribute at the same $\lambda$: it is confined mainly by the zero-point fluctuations of the gauge modes through the term $G^{ij}\Phi^\dagger M_iM_j\Phi$, which acts as a mass $\propto g_2^2\langle x^2\rangle$ for $r$. %
The fraction depends weakly on $\lambda$.
It is $0.2993$, $0.2999$ and $0.3043$ for $\lambda=0$, $0.01$ and $0.1$; at the Standard Model value $\lambda(\mub)=-0.0098$, where the Higgs zero mode is a narrow resonance (Secs.~\ref{rep:h-matching}--\ref{rep:h-reduced}), the real part of the resonance energy lowers $Q$ by about $0.08$ relative to $\lambda=0.01$, which gives $\fzm=0.298$. %

\subsection{Couplings}
\label{rA:zm-couplings}

The values above use one-loop couplings at $1/\Lb$, $g_2=0.5268$, $g_3=0.5396$ and $g_Y=0.4415$, with $Y_H=\tfrac12$; the top Yukawa coupling of this set is $y_t=0.4585$. %
The NNLO couplings of assumption~\ref{assump:A2} at $\mub$ are $g_2=0.5278$, $g_3=0.5357$ and $g_Y=0.4428$, within $0.8\%$ of these. %
At this order the zero-mode energies scale as $g^{2/3}$, so the difference moves $\fzm$ by less than $0.001$. %
Taking $\mu=2\pi/\Lb$ instead of $1/\Lb$ lowers $\fzm$ by $0.003$, and determining $\Lb$ self-consistently from $(1-\fzm)C$ changes it by $10^{-4}$. %

\subsection{Next order: the one-loop holonomy potential}
\label{rA:zm-nextorder}

At $\mathcal O(g^{4/3})$ the zero modes feel the curvature of the one-loop potential that the nonzero modes of charged fields create for constant connections~\cite{Luscher:1982ma,vanBaal:1988va} (Fig.~\ref{fig:rA-holonomy-loop}).
A constant connection along the generator $T^a$ shifts the twist of a charged field, $\bm\vartheta\to\bm\vartheta+\bm A\,q/2\pi$ for a charge $q$, so the curvature follows from the twist dependence of the image sum~\eqref{eq:rA-E1}.
Differentiating twice under the sum gives a factor $-(2\pi)^2n_in_j$, and at the cube the cubic average $\sum'n_in_jf(|\bm n|)=\tfrac13\delta_{ij}\sum'|\bm n|^2f(|\bm n|)$ turns $(\bm n^{\mathsf T}\bm n)^{-2}$ into $(\bm n^{\mathsf T}\bm n)^{-1}$.
A field with weight $d_f$, statistics $s_f$, twist $\bm\vartheta_f$ and ${\rm tr}_f\,T^aT^b=t_f\delta^{ab}$ ($t_f=Y^2$ for hypercharge) therefore contributes
\begin{equation}
L\frac{\partial^2E}{\partial A^a_i\partial A^b_j}\bigg|_0=\frac{\delta^{ab}\delta_{ij}}{3\pi^2}\sum_fd_fs_ft_f\,Z_{\rAone}(1;\bm\vartheta_f),
\label{eq:rA-holonomy-curvature}
\end{equation}
where $Z_G(s;\bm\vartheta)=\sum'\cos(2\pi\bm\vartheta\cdot\bm n)(\bm n^{\mathsf T}G\bm n)^{-s}$ is continued analytically to $s=1$.
Its values are $Z_{\rAone}(1;0)=-8.913633$, with the slow-mode cusp removed, and $Z_{\rAone}(1;\bm\vartheta_{\rm AAA})=-2.519356$. %
For periodic fields at trivial holonomy the charged mode with $\bm n+\bm\vartheta=0$ would make the potential non-analytic; that mode is part of the zero-mode Hamiltonian and is removed from the sum.
For $SU(3)$ with six Dirac flavours this reproduces van Baal's coefficient $-0.19631$~\cite{vanBaal:1988va}. %
For the Standard Model the curvatures are $-0.393$ for $SU(3)$, $-0.242$ for $SU(2)$, including the nonzero modes of the Higgs doublet, and $+0.700$ for hypercharge. %
The AAA fermions confine hypercharge, whose curvature is positive, while the periodic bosons dominate the negative non-abelian curvatures.

\begin{figure}[tbp]
\centering
\begin{tikzpicture}
\begin{feynman}
\vertex[crossed dot] (x1) at (-2.6,0) {};
\vertex (a) at (-1.0,0);
\vertex (b) at (1.0,0);
\vertex[crossed dot] (x2) at (2.6,0) {};
\diagram*{
(x1) -- [photon, edge label=\(A^a_i\)] (a),
(a) -- [fermion, half left, edge label=\(\bm n+\bm\vartheta_f\)] (b),
(b) -- [fermion, half left] (a),
(b) -- [photon, edge label=\(A^b_j\)] (x2),
};
\end{feynman}
\end{tikzpicture}
\caption{The one-loop curvature of the holonomy potential, Eq.~\eqref{eq:rA-holonomy-curvature}.
The loop runs over the nonzero modes of every charged field, with momenta $2\pi(\bm n+\bm\vartheta_f)$ on the torus; the crossed dots are insertions of the constant gauge zero modes.
Charged scalars and vectors contribute in addition through the seagull vertex with both insertions at one point.
Summed over the modes, the diagram is the second twist derivative of the image sum, which turns the $s=2$ Epstein function of the Casimir energy into the $s=1$ function.}
\label{fig:rA-holonomy-loop}
\end{figure}

\paragraph{First-order shift of the zero-mode energy.}
In first-order perturbation theory the curvature shifts the ground energy by $\Delta Q=\tfrac12\sum_{\rm G}M_{\rm G}\langle A^2\rangle_{\rm G}$, where $M_{\rm G}$ is the curvature per component and $\langle A^2\rangle_{\rm G}$ the summed spread of the ground state of factor G.
With the Gaussian spreads $\langle A^2\rangle=3.66$ for $SU(3)$ ($3.77$ from the virial theorem) and $1.40$ for $SU(2)$, and $\langle B^2\rangle=0.82$, the first-order shift is $\Delta Q=-0.72-0.17+0.29=-0.60$, or $\Delta\fzm=-0.014$; a virial spread for $SU(3)$ gives $-0.015$. %
The gluon and $W$ modes lower $\fzm$ by $0.021$, and hypercharge, which the AAA fermions confine, raises it by $0.007$. %
The dependence of the full one-loop potential on the Higgs radius $r$ enters the Higgs zero mode as cross terms (Sec.~\ref{rep:h-cross}).
The nonzero modes also give the Higgs zero mode a mass, which Sec.~\ref{rep:h-matching} treats together with the one-loop finite-volume potential along $r$.
In the Gaussian estimate used here this mass lowers $\fzm$ by a further $0.003$, and the NNLO top Yukawa coupling at $\mub$, $y_t=0.428$, changes that by $0.0003$. %
Together the first-order terms give $\fzm\simeq0.28$. %

\paragraph{Convergence of the small-volume expansion.}
The expansion parameter $g^{2/3}\simeq0.66$ is not small, but the published $SU(2)$ series converges quickly at this coupling. %
For three flavours its corrections are $-11\%$, $-0.4\%$ and $-0.15\%$ of the leading term at $g=0.53$. %
For $SU(3)$ with three flavours perturbation theory is expected to fail only near $g\simeq0.64$, for the first excited level~\cite{vanBaal:1988va}. %
We keep $\fzm=0.30$ as the benchmark and quote the first-order value $\fzm\simeq0.28$ as the next-order variant. %

\subsection{Shape dependence at next order}
\label{rA:zm-shape}

The next-order terms depend on the shape through twisted Epstein sums at $s=1$.
Section~\ref{rep:shape} shows that at fixed volume $Z_G(1;\bm\vartheta)$ is an eigenfunction of the shape Laplacian with eigenvalue $-\tfrac43$, so these terms enter the no-minimum argument in a controlled way.
Their Hessians at the cube follow from those of the twisted Epstein sums along the aspect and angle geodesics, $9.050$ and $-2.072$ for $Z_G(1;0)$ and $-8.425$ and $6.736$ for $Z_G(1;\bm\vartheta_{\rm AAA})$, whose traces are $-\tfrac43$ times the sums; the same code reproduces $\lambda_E$ and $\lambda_{T_2}$ of Table~\ref{tab:rA-spin} at $s=2$. %
The matching energy then has aspect and angle Hessians $7.75/L$ and $-4.90/L$, and the Higgs mass adds $0.60/L$ and $-0.34/L$. %
Subtracting $\pi^2L$ times these from the Hessian eigenvalues of $Z_{\rm eff}$ gives $\kappa_E=18.0$ and $\kappa_{T_2}=-8.34$ at $\fzm\simeq0.28$. %
The Hessian restores the aspect modes slightly: it lowers the aspect stiffness $\kappa_E$ of Sec.~\ref{rep:shape} by $4\%$, and together with the smaller zero-mode energy brings it from $19.32$ to $18.0$. %
The Higgs-mass part accounts for only $0.14$ of the shift of $\kappa_E$, and the virial spread for $SU(3)$ changes these results by $0.2\%$. %
The first-order estimate, $\fzm\simeq0.28$ together with this Hessian, lowers the shape fossil $\sigma_E$ of Sec.~\ref{rep:quantum-shape} by about $10\%$, by $6\%$ from the Hessian and $4\%$ from the energy, which we treat as a one-sided uncertainty at next order; higher orders are not estimated. %

\subsection{Consequences for the background and for the Higgs zero mode}
\label{rA:zm-consequences}

\paragraph{Effective weight and stiffnesses.}
At leading order $Q$ does not depend on the shape, so the zero modes leave the shape Hessian unchanged, lower the signed weight to
\begin{equation}
\rANeff=(1-\fzm)\NC=17.9 , %
\end{equation}
and raise every normalized shape stiffness $\kappa_A$ of Sec.~\ref{rep:shape} to $\kappa_A/(1-\fzm)$.
With $\fzm=0.2999$ the leading-order stiffnesses are $\kappa_E=19.324$ and $\kappa_{T_2}=-9.074$, with $\sum_A\kappa_A=8/(1-\fzm)=11.43$, and the shape commutator $\delta$ of Sec.~\ref{rep:quantum-shape} grows by $(1-\fzm)^{-3/4}=1.307$ relative to free fields. %
We quote results for free fields, $\fzm=0$, and with the zero-mode energy, $\fzm=0.30$; Table~\ref{tab:rA-fzero} of Sec.~\ref{rep:turning} collects the background quantities that change. %

\paragraph{Adiabatic following.}
The zero modes follow the shape adiabatically.
In conformal time the shape changes at a rate of order $h\Lb\sigma_E\simeq5\times10^{-4}$, far below their gaps, which are of order $g^{2/3}$ for the non-abelian modes and about $0.4$ for the softer hypercharge mode (harmonic estimate $0.36$, variational bound $0.46$). %
In the variables $y=g^{1/3}x$, in which the pure-gauge Hamiltonian is $g^{2/3}(\tfrac12\pi_y^2+\tfrac14F[y]^2)$, the copies of the vacuum under large gauge transformations lie at $|y|\simeq4\pi/g^{2/3}\simeq19$, far outside a ground state of rms width about $0.6$ per component. %
Tunneling to the copies is therefore negligible, and the global features that could make the zero-mode energy shape dependent (Sec.~\ref{rA:zm-isospectral}) do not act.
The excited levels are shape independent too, by Proposition~\ref{prop:rA-isospectral}; Sec.~\ref{rep:gate} returns to them.

\paragraph{The Higgs zero mode.}
The Higgs zero mode is the one constant mode whose fate this energy count does not settle.
At $\mub$ the Standard Model quartic is negative, $\lambda(\mub)=-0.0098$ (assumption~\ref{assump:A2}), so the Hamiltonian~\eqref{eq:rA-H0} has no normalizable ground state along the Higgs direction. %
The zero-point fluctuations of the electroweak gauge modes still confine the Higgs zero mode near the origin, and for $\lambda$ above a finite-volume threshold it is a narrow resonance.
Sections~\ref{rep:h-matching}--\ref{rep:h-universality} compute that threshold and the survival of the resonance through the turn.
A curvature coupling is negligible at the turn itself, where $R=6(\dot H+2H^2)=12h^2$ and $R\Lb^2=12(h\Lb)^2\simeq2\times10^{-4}$, but late in the expanding branch, $\tau\gtrsim4$, it decides whether the zero mode stays near the origin (Sec.~\ref{rep:h-universality}). %
